\chapter{\textbf{Introducción}}
\label{ch:Intro}
\thispagestyle{fancy}

\section{Contexto y motivación}
\label{sec:contexto_motivacion}

En los últimos años, la digitalización ha dejado de ser una opción reservada a grandes empresas y se ha convertido en una necesidad también para pequeños negocios. Sectores como la hostelería dependen cada vez más de herramientas digitales para gestionar reservas, responder consultas, organizar horarios y mejorar la atención al cliente. En este contexto, los restaurantes se enfrentan a un reto habitual: ofrecer una atención rápida y constante sin aumentar de forma excesiva la carga de trabajo del personal. Esta tendencia general de transformación digital puede situarse dentro de un marco más amplio de evolución del tejido productivo europeo y español, tal y como reflejan los informes recientes de la Comisión Europea y del Instituto Nacional de Estadística \cite{CE2024DigitalDecade,INE2025TIC}.

La atención telefónica sigue siendo un canal importante para muchos clientes, especialmente cuando quieren realizar una reserva, modificarla, consultar horarios o resolver dudas sencillas. Sin embargo, atender llamadas de forma manual puede interrumpir el trabajo diario del restaurante, sobre todo en horas de mayor actividad. Además, cuando el personal no puede responder a tiempo, se pueden perder reservas o generar una mala experiencia para el cliente.

La evolución reciente de la inteligencia artificial, especialmente de los modelos de lenguaje y de las tecnologías de voz, abre la posibilidad de crear asistentes conversacionales más naturales que los sistemas automáticos tradicionales. Estos sistemas permiten interpretar mejor lo que el usuario quiere decir, mantener cierto contexto durante la conversación y ejecutar acciones concretas, como consultar disponibilidad o registrar una reserva. Según informes recientes, el uso de inteligencia artificial en las organizaciones está creciendo de forma notable, aunque su integración real en los procesos de negocio todavía sigue siendo desigual \cite{McKinsey2025StateAI}.

Este proyecto nace de esa oportunidad: aplicar tecnologías actuales de inteligencia artificial conversacional a un caso práctico y cercano, como la gestión telefónica de reservas en un restaurante. La motivación principal no es sustituir por completo la atención humana, sino reducir tareas repetitivas, mejorar la disponibilidad del servicio y ofrecer una primera respuesta automática que pueda adaptarse a las necesidades del negocio. De esta forma, el sistema puede actuar como apoyo al personal del restaurante y como punto de partida para futuras integraciones con herramientas internas de gestión.

Por tanto, el interés del proyecto se encuentra en desarrollar una solución general, funcional y próxima a un entorno real, pero lo bastante flexible como para ser adaptada posteriormente a cada restaurante. Cada negocio puede tener sus propios horarios, normas de reserva, sistemas internos o forma de trabajar, por lo que el objetivo no es construir un producto cerrado e idéntico para todos los casos, sino una base software sólida sobre la que realizar esas adaptaciones.

\section{Problemática actual: Limitaciones de los IVR tradicionales ("Pulse 1...") y chatbots de texto}
\label{sec:problematica_actual}

\textoBase

\textoBaseAlt

\section{Objetivos del proyecto}
\label{sec:objetivos_proyecto}

\textoBase

\subsection{Objetivo general}
\label{subsec:objetivo_general}

\textoBaseAlt

\subsection{Objetivos específicos (técnicos y funcionales)}
\label{subsec:objetivos_especificos}

\textoBase

\section{Alcance del proyecto (Caso de uso: Restaurante)}
\label{sec:alcance_proyecto}

\textoBase

\textoBaseAlt

\section{Estructura del documento}
\label{sec:estructura_documento}

\textoBase